\section{Objetivos de la Propuesta}
 
\subsection{Objetivo General}
 
\begin{quote}
\textit{Rediseñar y optimizar el Sistema PLADDES de Trámite Documentario Virtual de la UNSAAC (\textit{tramite.unsaac.edu.pe}), migrando su arquitectura tecnológica a un stack moderno y mejorando la experiencia de usuario, la integración con el catálogo TUPA, el seguimiento de expedientes y la gestión administrativa; mediante el desarrollo de una plataforma web transaccional que automatice los procedimientos establecidos en el TUPA institucional, mejorando la eficiencia, transparencia y seguimiento de los trámites universitarios.}
\end{quote}
 
\subsection{Objetivos Específicos}
 
\begin{enumerate}
    \item \textbf{Migrar el frontend} de AdminLTE 2/Bootstrap 3/jQuery a \textbf{React.js + Tailwind CSS}, eliminando la deuda tecnológica y mejorando el rendimiento y mantenibilidad del sistema. \textit{(Atiende P1)}
 
    \item \textbf{Rediseñar el flujo de registro de solicitudes} con validación en tiempo real, mensajes de error claros y confirmación de envío explícita para el usuario. \textit{(Atiende P2)}
 
    \item \textbf{Implementar un panel de seguimiento de expedientes} con vista de línea de tiempo, historial de estados y dependencia responsable visible en cada etapa. \textit{(Atiende P3)}
 
    \item \textbf{Desarrollar un módulo de verificación de pagos} que permita al sistema cruzar el código de pago con los registros de caja, reduciendo la validación manual. \textit{(Atiende P4)}
 
    \item \textbf{Implementar notificaciones automáticas por correo electrónico} al registrar una solicitud y en cada cambio de estado del expediente. \textit{(Atiende P5)}
 
    \item \textbf{Establecer un mecanismo formal de gestión de mantenimientos programados} con avisos anticipados y página de estado del servicio. \textit{(Atiende P6)}
 
    \item \textbf{Integrar el catálogo digital del TUPA} (procedimientos, requisitos, costos y plazos) directamente en el flujo de registro, de forma contextual según el tipo de trámite seleccionado. \textit{(Atiende P7)}
 
    \item \textbf{Implementar guía de documentos requeridos} por tipo de trámite, indicando al usuario exactamente qué archivos adjuntar antes de enviar su solicitud. \textit{(Atiende P8)}
 
    \item \textbf{Desarrollar un panel administrativo} para las dependencias (Decanatos, VRAC, DIGA, UTH, etc.) con bandeja de expedientes, asignación de responsables, reportes y estadísticas de productividad. \textit{(Atiende P9)}
\end{enumerate}